\chapter{Revisão da Literatura}

A revisão da literatura tem como finalidade contextualizar os principais conceitos relacionados ao desenvolvimento do sistema proposto para a loja Tropical Mix, abordando a automação de processos comerciais, o controle de estoque, o uso de sistemas web e a digitalização de operações em pequenos e médios empreendimentos. Dessa forma, busca-se apresentar o embasamento teórico necessário para compreender o problema identificado e justificar a adoção de uma solução tecnológica voltada à gestão de estoque e vendas.

\section{Histórico da Gestão Automatizada}

A informatização de processos empresariais teve início com a popularização dos sistemas de gestão nos anos 1980, quando grandes corporações passaram a utilizar softwares para controle de produção e finanças, conhecidos como \textit{Enterprise Resource Planning} (ERP). Com o passar dos anos, esses sistemas evoluíram e começaram a atender também pequenas e médias empresas, proporcionando maior controle sobre estoques, vendas e relacionamento com clientes \cite{laudon2022}.

Inicialmente, os sistemas eram instalados localmente (\textit{on-premise}), exigindo infraestrutura e manutenção interna. Entretanto, com o avanço da computação em nuvem e a consolidação de serviços baseados na \textit{web}, emergiu uma geração de soluções que combinam acesso remoto — \textit{anytime, anywhere} — com a centralização do armazenamento e da gestão de dados na infraestrutura do provedor \cite{armbrust2009}. Em linha com essa visão, a definição do NIST caracteriza a nuvem como um modelo que habilita acesso ubíquo e sob demanda, via rede, a um conjunto compartilhado de recursos computacionais \cite{mell2011}.

No contexto do comércio varejista, a adoção de sistemas de controle de estoque e vendas passou a representar não apenas um diferencial competitivo, mas uma necessidade operacional. Tais ferramentas garantem precisão nas informações, redução de perdas e agilidade na tomada de decisões, fatores essenciais para o crescimento de uma empresa.

\section{Gestão na Atualidade}

A gestão empresarial na atualidade está cada vez mais associada ao uso estratégico da tecnologia para otimizar processos e garantir maior eficiência operacional. Em um cenário de alta competitividade, empresas de pequeno e médio porte, como a \textbf{Tropical Mix}, precisam adotar soluções tecnológicas que auxiliem no controle de informações, reduzam falhas humanas e aumentem a produtividade.

A utilização de sistemas web tem se mostrado uma alternativa acessível e eficaz para modernizar o gerenciamento de negócios locais. Esses sistemas permitem a automatização de tarefas antes realizadas manualmente, como o controle de estoque e o registro de vendas, além de oferecerem acesso remoto e informações centralizadas em tempo real. Segundo \cite{mell2011}, a computação em nuvem possibilita que dados e aplicações sejam acessados de qualquer lugar, garantindo flexibilidade, segurança e integração entre diferentes setores da empresa.

De forma semelhante, \cite[p.~2]{armbrust2009} destacam que a computação em nuvem permite que aplicações empresariais sejam executadas e gerenciadas de forma centralizada, eliminando a necessidade de infraestrutura local complexa. Para empresas como a \textbf{Tropical Mix}, que antes realizavam o controle de estoque e vendas manualmente, essa abordagem representa uma oportunidade de reduzir custos e melhorar a precisão das informações.

Além disso, estudos da Organização para a Cooperação e Desenvolvimento Econômico (OCDE) ressaltam que a digitalização dos processos empresariais aumenta a competitividade e a capacidade de inovação das organizações \cite{oecd2014}. A centralização de dados e o acesso em tempo real permitem que os gestores tomem decisões mais rápidas e embasadas, o que é essencial para a sustentabilidade de micro e pequenas empresas.

Dessa forma, observa-se que a tecnologia da informação deixou de ser apenas uma ferramenta de apoio e passou a ocupar papel central na administração de pequenas empresas. O sistema proposto para a \textbf{Tropical Mix} busca justamente incorporar esses avanços, oferecendo uma plataforma intuitiva, automatizada e acessível, capaz de transformar a gestão operacional em um processo mais ágil, confiável e orientado a dados.


\section{Gerenciamento em Diversos Setores}

A tecnologia da informação tornou-se um elemento essencial para o gerenciamento de empresas em diversos setores, permitindo maior integração entre processos e acesso rápido a informações estratégicas. No comércio varejista, em especial, a automação das operações tem proporcionado melhor controle de estoque, registro de vendas e acompanhamento de desempenho financeiro. Segundo \cite{obrien2013}, os sistemas de informação são fundamentais para apoiar a tomada de decisões e garantir a eficiência das atividades empresariais, independentemente do porte da organização.

De forma semelhante, \cite{laudon2022} destacam que a utilização de sistemas web e plataformas digitais possibilita que empresas otimizem suas rotinas administrativas e mantenham seus dados centralizados, promovendo maior agilidade e precisão nas informações. No caso da loja \textbf{Tropical Mix}, a implementação de um sistema web para controle de estoque e vendas representa a aplicação prática dessas tecnologias, permitindo substituir o controle manual por uma ferramenta automatizada, acessível e voltada à melhoria contínua da gestão empresarial.


